% =====================================================================
% Appendix: Preregistration detail
% -- Relocated from sections/preregistration.tex on 2026-07-26 as part of the
%    structural trim. NOTHING WAS DELETED IN THE MOVE: the disclosed-contact
%    subsection, all seven interpretive choices, and the full spot-check
%    narrative appear here verbatim, with their % SOURCE comments.
%
% PRIMARY SOURCES:
%   docs/AUDIT_REGISTRATION_2026-07-15.md     (frozen; Amendment 1 appended)
%   docs/ATLAS_MINING_REGISTRATION_2026-07-15.md (frozen)
%   docs/AUDIT_VERDICTS_2026-07-20.md §"Interpretive choices", §Provenance
%   docs/IDENTICAL_SCORE_CHURN_2026-07-21.md §"Why this note exists"
%   docs/ATLAS_REV2_CORRECTION_2026-07-21.md  (spot-check, rev-1 -> rev-2 delta)
%
% DISCLOSURE BOUNDARY -- ENFORCE ON EVERY EDIT.
% The retired "6.3% flips at identical score" anecdote (bnb-4bit same-repo
% rerun) may appear in THIS FILE ONLY, labelled as (i) pre-registration data
% contact and (ii) outside the frozen 59-pair manifest. It must appear nowhere
% else in the paper -- including sections/preregistration.tex, which held it
% until 2026-07-26 -- and nowhere at all in the blog post. It is a motivating
% anecdote inside the honest-disclosure narrative, never an atlas statistic.
% THE BOUNDARY TRAVELS WITH THE ANECDOTE: if this paragraph is ever moved again,
% move this comment with it and update the pointer in preregistration.tex.
% =====================================================================

\section{Preregistration detail}
\label{app:prereg-detail}

% ITEM COUNT UPDATED 2026-08-13 with the two relocations of that date
% (app:prereg:whatfixed, and the body's account of choice 1 inside
% app:prereg:choices). Six items became seven.
This appendix holds the full text of seven items summarised in
\S\ref{sec:prereg}: the freeze timeline, what the freezes fix in advance, the
disclosed pre-registration data contact, the seven interpretive choices, the
result-inspection discipline, the H3 reporting rule as registered, and the
spot-check narrative.

\subsection{The freeze timeline}
\label{app:prereg:timeline}

% RELOCATED 2026-08-04 (flagship narrative, operation 4) from
% sections/preregistration.tex. The table and its caption are moved verbatim;
% nothing was deleted, and the body keeps a pointer to it.
% SOURCE: PREREGISTRATION.md header lines 3-7;
% docs/RESULTS_2026-07-15_ATLAS_AUDIT.md §3 (commits b74fd58, d6e02dd, f06348f,
% 715a7ce); docs/AUDIT_VERDICTS_2026-07-20.md §Provenance (claim-table sha256).
\begin{table}[!t]
\centering
\small
\caption{Freeze timeline. Each artifact was committed before the analysis it
governs could be run, and each frozen file carries a \emph{Dated Amendments}
section: deviations are appended, never edited in, and each records whether
results had been inspected before the decision.}
\label{tab:freeze-timeline}
% WIDTH FIX 2026-08-05: tightened column separation and wrapped the widest
% text column. No cell value changed and no row was dropped.
\setlength{\tabcolsep}{4pt}
\small
\begin{tabular}{lll>{\raggedright\arraybackslash}p{4.6cm}}
\toprule
Date & Artifact & Commit & What it fixed in advance \\
\midrule
2026-07-11 & \texttt{PREREGISTRATION.md} & --- & Grid, metrics, TOST margin (2\,pp), H3 decision rule \\
2026-07-15 & Three registrations & \texttt{b74fd58} & Audit, atlas-mining, and mini-grid protocols \\
2026-07-15 & Audit Amendment 1 & \texttt{d6e02dd} & Blind-extraction independence mechanism \\
2026-07-15 & Pair manifest & \texttt{f06348f} & The 59 atlas pairs, before any flip statistic \\
2026-07-15 & Claim table & \texttt{715a7ce} & The 17 audited claims, before any verdict \\
2026-07-20 & WikiText-2 amendment & --- & ``Document'' definition, Decision Point A \\
\bottomrule
\end{tabular}
\end{table}

\subsection{What the freezes fix in advance}
\label{app:prereg:whatfixed}

% RELOCATED 2026-08-13 from sections/preregistration.tex
% §"What was frozen, and when" (\S\ref{sec:prereg}), under the reviewer's
% structural-streamlining request. Moved VERBATIM; nothing was deleted, and the
% body keeps a two-sentence compliance statement plus a pointer here.
% TWO MECHANICAL EDITS, BOTH RECORDED: the source paragraph's cross-reference
% read "Table~\ref{tab:freeze-timeline} in Appendix~\ref{app:prereg-detail}",
% which is repointed to the local table now that the text sits inside that
% appendix; and its closing pointer sentence ("The appendix carries the rest:
% ...") is dropped, because it pointed AT this appendix from the body and the
% body now carries that pointer. Its opening sentence ("Each artifact ... before
% the decision") is not repeated here: it is reproduced word for word in the
% caption of Table~\ref{tab:freeze-timeline} directly above, and a condensed
% form of it is what the body retains.
% SOURCE: PREREGISTRATION.md header lines 3-7;
% docs/RESULTS_2026-07-15_ATLAS_AUDIT.md §3 (commits b74fd58, d6e02dd, f06348f,
% 715a7ce); docs/AUDIT_VERDICTS_2026-07-20.md §Provenance (claim-table sha256).
Table~\ref{tab:freeze-timeline} dates every freeze against its commit. The
frozen claim table's content hash is recorded with the verdicts (sha256
\texttt{842b9756...5af7b15}), as is the atlas summary's
(\texttt{98201ade...10a4712d}) and the container image's, so a reader can
confirm that the inputs to any reported number were the frozen ones.

Three properties of that arrangement do work in the paper. The mechanical parts
of the audit, meaning the inclusion trigger vocabulary, the extraction fields,
the verdict formulas and the robustness sweep, were fixed before any claim's
power was computed, so the audit cannot have selected its claims or its
statistics to produce a headline. The atlas pair list was fixed before any flip
statistic, so the atlas cannot have selected pairs to produce churn. And the H3
decision rule was fixed, in algebra, before the first compressed checkpoint
existed, so the paper reports whatever that rule returns.

\subsection{Result-inspection discipline}
\label{app:prereg:discipline}

% RELOCATED 2026-08-04 (operation 4) from sections/preregistration.tex §3.4,
% moved verbatim.
% SOURCE: CLAUDE.md/AGENTS.md project guardrails; PREREGISTRATION.md
% §"Outcomes and Analysis"; docs/MINIGRID_REGISTRATION_2026-07-15.md §5.
During grid execution only job health, checksums, expected-file coverage and
receipt pairing are inspected; the first accuracy inspection happens only after
the mini-grid validator passes over the complete expected file set. The
registered analysis is run once per cell, the escalation rule applied
mechanically, and a dated escalation decision record written the same day. The
paired-bootstrap rank-flip denominator convention (tie replicates included,
ties also reported separately) was fixed in words in the registration, naming
the implementing commit, so it cannot be re-chosen after inspection. No
registered analysis is tuned after results are seen: not the calibration
builder, not the paired bootstrap, not the verdict rule.

\subsection{The H3 reporting rule, stated before the results exist}
\label{app:prereg:h3rule}

% RELOCATED 2026-08-04 (operation 4) from sections/preregistration.tex §3.5,
% moved verbatim except for one repointed cross-reference
% (sec:minigrid:escalation -> app:escalation).
% SOURCE: PREREGISTRATION.md §"H3 Decision Rule";
% docs/MINIGRID_REGISTRATION_2026-07-15.md §§1, 3, 4.
% THIS SUBSECTION STATES NO H3 OUTCOME. That is deliberate and survives the
% signed verdict of 2026-07-26: this is the protocol, §\ref{sec:minigrid} is the
% result.
The registered decision rule is defined over eight confirmatory
model-by-benchmark cells:
$\{$Qwen2.5-1.5B, Qwen2.5-7B, Llama-3.2-3B, Llama-3.1-8B$\}
\times \{$MMLU, GSM8K$\}$ at 4 bits, comparing GPTQ seed $s$ against AWQ
seed $s$ on byte-identical calibration samples. H3 is \emph{supported} if winner
flips occur in at least 3 of the 8 cells, or the range/gap criterion holds in at
least 4 of the 8. It is \emph{disconfirmed} if winner flips occur in at most 1
of the 8 and $\max(\mathrm{range}_{\mathrm{GPTQ}}, \mathrm{range}_{\mathrm{AWQ}})
< 0.5 \times \mathrm{gap}$ in at least 6 of the 8. Everything else is
\emph{inconclusive}, reported without post-hoc promotion.

The mini-grid executed \textbf{4 of those 8 cells}
($\{$Qwen2.5-1.5B, Llama-3.2-3B$\} \times \{$MMLU, GSM8K$\}$), with the 7B/8B
cells deferred behind a pre-committed, mechanical escalation rule. The mini-grid
registration stated the consequence of each branch in advance: if all eight
cells complete, the frozen rule is applied exactly as registered; if they do
not, the paper reports the four completed cells descriptively and H3 is
undecided under the registered rule, never supported or disconfirmed on four
cells, and no reduced-cell variant of the rule is constructed after results are
seen.

The escalation screen fired on 2026-07-23, the 7B/8B cells were built, and
\textbf{all eight cells completed}, so the second branch never applied and the
frozen rule was applied once over the full set. Both decisions are dated, signed
records
(\texttt{docs/MINIGRID\_ESCALATION\_DECISION\_2026-07-23.md},
\texttt{docs/H3\_EIGHT\_CELL\_DECISION\_2026-07-26.md}), and the screen's own
authorship discipline is worth stating: it decides which cells to build and
states no H3 outcome, so the confirmatory rule was never applied to a cell set
selected after its result was known. The screen's per-cell outcomes are in
Appendix~\ref{app:escalation} and the verdict in
\S\ref{sec:minigrid:verdict}; this subsection states neither.

\subsection{Disclosed pre-registration data contact}
\label{app:prereg:contact}

Preregistration is a claim about ordering, so the places where our order was
imperfect are disclosed by the registrations themselves rather than discovered
by a reader.

% SOURCE: docs/AUDIT_REGISTRATION_2026-07-15.md §1.
\textbf{Audit.} Five candidate claims were collected with exact quotes on
2026-07-15 during a feasibility sweep, before the registration was written
(GPTQ, LLM.int8(), and SmoothQuant abstracts; the Red~Hat AI W4A16 Llama-3.1-8B
model card; Meta's quantized-Llama blog). No power computation had been run on
any of them. They enter the pool through the same §3 criteria as
later-collected claims, and they are not treated as a separate stratum, because
the contact was with the \emph{quotes}, not with any verdict quantity.

% SOURCE: docs/ATLAS_MINING_REGISTRATION_2026-07-15.md §1; the 99 probe cells
% are excluded per §6 and per docs/AUDIT_VERDICTS_2026-07-20.md ruling #6.
\textbf{Atlas.} Two feasibility probes were run on 2026-07-15 before the
registration was drafted, and their results were known: TheBloke/Llama-2-7B-GPTQ
against meta-llama/Llama-2-7b-hf on ARC-Challenge (74 flips of 1{,}170, net
$-1.03$\,pp), and a Neural~Magic Llama-3.1-8B baseline against its W4A16
quantization on \texttt{bbh\_boolean\_expressions} (17 flips of 250, net
$+1.2$\,pp). Those pairs are flagged in the atlas and excluded from every
headline aggregate; the exclusion covers all 99 cells belonging to the two probe
\emph{pairs}, not merely the two probe tasks (\S\ref{app:prereg:choices},
interpretive choice~6).

% SOURCE: docs/IDENTICAL_SCORE_CHURN_2026-07-21.md §"Why this note exists".
% DISCLOSURE BOUNDARY: this paragraph is the ONLY place the 6.3% anecdote may
% appear in any FlipEval public-facing text. Never in the blog, never in the
% abstract, never in §\ref{sec:atlas}, never in §\ref{sec:prereg}.
\textbf{A retired anecdote.} The project's original motivating observation was a
feasibility rerun of a bnb-4bit checkpoint against its own repository's
full-precision weights, in which roughly 6.3\% of items flipped correctness while
the aggregate score was unchanged. That observation is what made the project
seem worth doing, and it is reported here for that reason. It is
\textbf{retired from every quantitative use}, on two independent grounds: it is
pre-registration data contact, and it lies outside the frozen 59-pair manifest,
since same-repository multi-precision runs are an amendment candidate that was
never acted upon. The registered replacement is the identical-score statistic of
\S\ref{sec:atlas:identical}, computed on the frozen atlas population: 145 of
1{,}707 cells (8.49\%) at a median churn of 7.20\%. Where an anecdote and a
registered statistic point the same way, the paper cites the statistic.
% REV-1 SURVIVOR FLAGGED 2026-07-26, RESOLVED 2026-07-26 on the author's ruling.
% All four figures in the preceding sentence were rev-1
% (results/identical_score_churn.csv: analysable_cells 1155, zero_delta_cells
% 113, zero_delta_share 0.097835, churn_median 0.062176) and are now rev-2
% (results/identical_score_churn_rev2.csv: 1707, 145, 0.084944, 0.072000).
% This passage cites the registered replacement statistic as CURRENT, not as
% history, so it must track rev-2 -- unlike \S\ref{sec:atlas:coverage}'s citation
% of the retracted 643, which is deliberate historical narration and stays.
% The abstract's and introduction's companion figure "a median of 6.2%" was the
% same rev-1 churn median and moved to 7.2% in the same commit; all three now
% state churn_median 0.072000 at their respective precisions.

\subsection{The interpretive choices that moved the headline}
\label{app:prereg:choices}

Frozen protocols do not eliminate analyst judgement; they make it visible and
datable. Seven passages of the audit registration were ambiguous enough to
require a ruling before verdicts could be computed. All seven were ruled on
2026-07-20, each is implemented in code, and each is reversible by re-running
with the alternative. Four matter for how the results should be read.

\paragraph{Choice 1: which margin is ``the applicable margin'', and the
$K = 1 \to 5 \to 4$ correction.}
% SOURCE: docs/AUDIT_VERDICTS_2026-07-20.md §"Interpretive choices" #1 (the
% SUPERSEDED record); resolution in docs/AUDIT_REGISTRATION_2026-07-15.md
% "Dated Amendments" Part 1, signed 2026-07-31; values in
% results/audit_verdicts_rev3.csv.
%
% THIS CHOICE IS RETIRED. Amendment 2 resolved it, and this paragraph is now a
% record of a resolved ambiguity rather than a live interpretive option. The
% previous version argued FOR the own-margin reading in the present tense,
% including the "0.15 pp" sentence the advisor review identified as the
% load-bearing error. That sentence is retained ONLY as the quoted reasoning
% that was found wrong; it must never be restated as the paper's position.
The frozen §4 names the 2\,pp registered margin first, adds ``(and at the
claim's own margin when it states one)'', and then labels the verdict
``underpowered for its own assertion'' \emph{at the applicable margin}. We read
that as leaving two live options, and chose the second. That was an error, and
Amendment~2 corrects it.

The conditional in the frozen text is the whole of it: \emph{when it states
one}. No audited source states one (\S\ref{sec:audit:taxonomy}). What the
superseded implementation substituted was each source's largest reported
$|\Delta|$, and the reasoning it ran on was that ``a source asserting parity
within 0.15\,pp has made a 0.15\,pp claim''. It has not. A reported delta is an
outcome of an evaluation that has already been run; a margin is a tolerance
fixed before it runs, and a claim cannot be held to a decision rule its authors
never adopted. Treating the one as the other inverted the direction of the
inference and, because a measured delta is typically far tighter than any
tolerance a author would have declared, it also made the test far harsher than
the registration authorised.

% The sequence of values is the evidence. Reported plainly, with no argument
% about what the analyst believed at each step (advisor review 3.3).
The sequence is reported exactly, because the order in which the numbers arrived
is itself the disclosure. The first pass applied the 2\,pp margin uniformly and
returned $K = 1$ of 12. Re-reading the frozen label produced $K = 5$; ruling on
choice~3 below moved one claim out of the determinate set, giving the $K = 4$
that was reported for ten days. Amendment~2 returns the analysis to the uniform
registered margin, where the rev-3 recomputation gives \textbf{1 of 11
assessable claims below the approximate planning threshold}, the same claim the
very first pass had flagged.

\paragraph{Choice 1, as the main text stated it.}
% RELOCATED 2026-08-13 from sections/preregistration.tex §"The interpretive
% choice that moved the headline" (\S\ref{sec:prereg:choices}), under the
% reviewer's structural-streamlining request. Moved VERBATIM, whole; nothing
% was deleted. The body keeps a three-sentence statement of the sequence and of
% Amendment 2's correction, because that sequence is the paper's canonical
% disclosure of an analyst decision that moved the headline and does not live
% only in an appendix; \S\ref{sec:audit:r04} states it a second time in the
% audit's own terms.
% IT OVERLAPS THE TWO PARAGRAPHS ABOVE BY CONSTRUCTION -- the body passage was a
% condensed retelling of them -- and is preserved whole rather than merged into
% them, so that the relocation loses no wording.
% SOURCE: docs/AUDIT_VERDICTS_2026-07-20.md §"Interpretive choices" #1 (the
% SUPERSEDED record); resolved by docs/AUDIT_REGISTRATION_2026-07-15.md "Dated
% Amendments" Part 1, signed 2026-07-31. Current values from
% results/audit_verdicts_rev3.csv.
%
% RESOLVED, NOT OPEN. The reasoning "a source asserting parity within 0.15 pp
% has made a 0.15 pp claim" is the error Amendment 2 exists to correct: R17's
% 0.15 is an OBSERVED DELTA (68.69 vs 68.54), not a declared tolerance. It must
% not be restated as the paper's position anywhere, in any file, in any wording.
The conditional decides it. \emph{When it states one}: no audited source states
one, so the applicable margin is the registered 2\,pp throughout, and what the
superseded implementation had substituted was each source's largest reported
delta. That substitution treats a measured outcome as a decision rule the source
never adopted. The sequence of values is reported plainly because it is the
disclosure: the first pass returned $K = 1$ of 12, the claim-specific reading
raised it to $K = 5$, the metric-compatibility ruling on R04 gave the $K = 4$
that was published, and the rev-3 recomputation at the registered margin returns
\textbf{\AuditBelowThresholdAtMedian{} of \AuditAssessable{} assessable claims
below the approximate planning threshold}, which is the claim the first pass had
flagged. All readings ship in the released CSV (\S\ref{sec:audit:results}).

\paragraph{Choice 2: what counted as the claim's own margin (withdrawn).}
% SOURCE: docs/AUDIT_VERDICTS_2026-07-20.md §"Interpretive choices" #2 (the
% SUPERSEDED record). Retained as the construction Amendment 2 withdraws.
%
% THE 2.0x-12.9x SHORTFALL RANGE IS WITHDRAWN HERE, NOT RECOMPUTED. It was
% required-n over reported-n AT THE RESULT-DERIVED MARGINS, so recomputing it at
% 2 pp does not produce a corrected version of the same quantity: ten of the
% eleven assessable claims have no shortfall at all, so there is no range to
% state. Do not reintroduce the sentence with new numbers in it.
This choice is withdrawn with the margins it operated on, and is recorded
because the withdrawn quantities were reported publicly.

Most claims cover several benchmarks with different deltas, so the superseded
implementation defined the claimed margin as the largest $|\Delta|$ the source
reports. Within its own frame that was the reading most favourable to the
source, granting the claim the widest margin its sentence could bear and
therefore the smallest required sample size. The point Amendment~2 makes is that
the frame was wrong to begin with: no choice among a source's reported deltas
yields a margin, because none of them is one.

Two published quantities fall with it. The shortfall range of $2.0\times$ to
$12.9\times$ was required $n$ over reported $n$ at those derived margins, and it
is \textbf{withdrawn rather than recomputed}; at the registered 2\,pp margin ten
of the eleven assessable claims have no shortfall at all, so there is no range
to restate. The MDD-to-margin ratios are likewise recomputed against the
registered margin in Table~\ref{tab:audit-mdd}, where the values themselves are
unchanged and only the denominators move.

What survives the withdrawal is the disposition of the remaining choices, which
were and remain resolved in the source's favour: pairing rather than independent
binomial variance (the paired assumption is the generous one, and
Table~\ref{tab:audit-mdd} shows the independent bound is uniformly worse), and
imputation by \emph{median} discordance from the most specific matching atlas
tier rather than by any upper quantile.

\paragraph{Choice 3: R04 is indeterminate rather than scored on a substituted
benchmark.}
% SOURCE: docs/AUDIT_VERDICTS_2026-07-20.md §"Interpretive choices" #3 and the
% R04 bullet in §Indeterminate; results/audit_verdicts_rev3.csv row R04.
% The K and shortfall figures below are PRE-AMENDMENT-2 HISTORY, quoted as what
% was decided at the time. Do not update them to rev-3 values.
Discussed in full at \S\ref{sec:audit:r04}, and unaffected by Amendment~2: R04
sits outside the registered binary paired-outcome calculation whichever margin
is applied. It is recorded here as a ruling because it is the second correction
that cost the paper a number. Under the rules then current it removed what the
first pass had reported as the audit's largest shortfall ($38.3\times$) and
moved the headline from $K = 5, J = 4$ to $K = 4, J = 5$; those figures are
quoted as history and are not current results. The GSM8K computation is retained
in the released CSV as a labelled transparency column so that the set-aside
quantity is inspectable rather than deleted.

\paragraph{Choice 5: TOST uses the one-sided $z$.}
% SOURCE: scripts/audit_stats.py::required_n_for_tost (Z_ONE_SIDED) and
% ::minimum_detectable_delta (two-sided), which are the implementation of
% record; \S\ref{sec:cert:method}. The rationale was first written up in
% docs/AUDIT_VERDICTS_2026-07-20.md §"Interpretive choices" #5 and
% docs/CERTIFICATION_TABLES_2026-07-20.md §Method, both REV-1 narratives whose
% numeric provenance is superseded; the z-choice argument they make is not
% population-dependent and is unaffected, but they are not the source for any
% number here.
The project's own helper, \texttt{flipeval.required\_n\_for\_effect}, uses the
two-sided $z_{1-\alpha/2}$. That is correct for \emph{detection} and wrong for
TOST, which rejects two one-sided nulls at level $\alpha$ each. Reusing the
helper unchanged would have inflated every required $n$ in this paper by about
27\%, in the conservative direction, under the name ``TOST'', and with no
symptom visible in any output. It would have made every audit shortfall and
every certification count larger than the correct value. The TOST formula is
therefore implemented locally with the one-sided $z_{1-\alpha} = 1.6449$,
documented, and cross-checked against the project implementation on the
quantities the two share; the paired standard deviation is pinned against
\texttt{flipeval.core.minimum\_detectable\_difference} on synthetic delta vectors
in the test suite, so the audit cannot silently fork from the tested library.

\paragraph{The remaining three rulings, in brief.}
% SOURCE: docs/AUDIT_VERDICTS_2026-07-20.md §"Interpretive choices" #4, #6, #7.
Choice~4: verdicts are computed at claim level rather than
claim\,$\times$\,benchmark, because the frozen table stores one pooled $n$ per
claim and operating at a finer granularity would mean inventing rows the freeze
does not contain; every summed task is listed in the released
\texttt{n\_basis} column. Choice~6: ``the two probe-tagged cells'' is read as
two probe \emph{pairs}, comprising 99 cells, all excluded. Choice~7: degenerate
rows are handled per input rather than discarded whole, so an indeterminate
claim keeps the components its available inputs support, reported as
supplementary and never verdict-bearing.

\subsection{Atlas validation and population correction}
\label{app:prereg-spotcheck}
% SOURCE: docs/ATLAS_REV2_CORRECTION_2026-07-21.md; the spot-check report of
% 2026-07-21 (10 cells, 262/262 fields); ruling R7.

The protocol required an independent spot-check of the atlas pipeline before any
atlas number could be quoted externally. We report what it found, because the
finding is the paper's thesis applied to the paper.

The check re-downloaded the raw per-item files for ten cells, stratified across
both sources and across zero-delta, high-churn, McNemar-significant and excluded
strata, and recomputed the joins, pairing gates, flip counts, churn, net deltas
and exact McNemar $p$ values from a \emph{fresh reimplementation of the
registered definitions}, deliberately not a rerun of our own pipeline, since
rerunning the same code confirms determinism rather than correctness. All
262 compared fields reconciled exactly, including dyadic $p$ values reproduced
without \texttt{scipy}, and no upstream data drift was detected.

The arithmetic was right. The \emph{population} was not. Two defects surfaced.
First, our implementation omitted a clause of our own registration: when the
prompt-hash gate fails, the protocol requires trying earlier run combinations in
reverse-chronological order, and the code simply took the newest run and gave
up. Eleven pairs contributed zero cells as a result. Crucially, the loss was
\textbf{not random}: it removed exactly those pairs whose quantized side had
been re-evaluated later, a property correlated with a model being popular enough
to re-run. Second, a newer details schema nested its metrics one level deeper
than our parser looked, so 583 cells were recorded as having no binary
correctness metric when in fact they had one, and a sentence in our own results
note attributed that gap to the upstream leaderboard's reporting standards. We
retracted the sentence.

We draw three points from this, and none of them is that we were unlucky.

\textbf{An aggregate can be exactly right and still be built on the wrong
population.} Every cell we checked was computed correctly. Had we validated by
recomputing our own numbers, which is the natural and useless form of
self-check, we would have confirmed all of them and shipped the selection bias intact. What
caught it was reconstructing the measurement from the protocol text rather than
from the code.

Preregistration did the work it is supposed to do. The registration
text was frozen before any statistic existed, and it fully determined the
direction of the repair: we did not choose a rule that flattered a result, we
executed a rule that was already binding and had been under-implemented. That
is why the fix is a correction rather than a post-hoc adjustment, and it is why
we can say so without asking to be trusted. The correction was nevertheless made
\emph{after} results had been inspected, and we disclose that plainly rather
than presenting the repair as a pre-specified step.

Both revisions are public. We publish rev-1 and rev-2 and report the
delta between them, rather than replacing the record with its corrected version.
A field in which near-lossless claims cannot be rechecked because per-item
outputs are never released, which is the finding of Section~\ref{sec:audit}, is
a field whose corrections are invisible. Ours is not, and the difference is the point of
the artifact rather than an apology for it.

\subsection{The rev-1 to rev-2 delta}
\label{app:prereg:rev2delta}
% SOURCE: docs/ATLAS_REV2_CORRECTION_2026-07-21.md §8 "Rev-1 -> rev-2 delta
% record", all four sub-tables (Population; Headline descriptives per stratum;
% Audit verdicts; Certification tables). Rev-2 ran as embers job 11341992
% (exit 0, 59/59 pairs), downstream regeneration as job 11343383 (exit 0).
% \revtwoTODO CLOSED 2026-07-26: this subsection is the narrative the marker
% asked for. Its own note recorded that every input already existed in §8 and
% that it was not blocked on any pending result.

The repair recovered cells rather than removing them, and the recovery was
confined to S1. The enumerated population is unchanged at 2{,}055 pair-task
cells; analysed cells rise from 1{,}254 to \textbf{1{,}807}, and the
probe-excluded analysable population from 1{,}155 to \textbf{1{,}707}. All of
that movement is S1 (846 $\to$ 1{,}398); S2 is unchanged at 309 cells, which is
the expected control, since no S2 cell's eligibility was affected by either
defect. The reverse-chronological run fallback was exercised for 1{,}007 cells
across 19 pairs, recorded per cell with the accepted run timestamps.

The descriptives moved in the direction a recovered population predicts rather
than in the direction that would flatter the paper. S1 median churn rises
slightly (0.1327 $\to$ 0.1375) and S1 median $|$net delta$|$ rises
(0.0226 $\to$ 0.0263), while the share of S1 cells certifiable at 2\,pp
\emph{falls} from 5.6\% to 4.9\% and the share showing a detectable difference
rises sharply from 17.5\% to 26.5\%, because the recovered pairs are larger-$n$
cells, which resolve more differences. Every S2 field is identical across revisions.

% HISTORY, PRE-AMENDMENT-2, AND SUPERSEDED. This paragraph reports what the
% rev-1 -> rev-2 population change did to the headline AS IT THEN STOOD. Both
% the count and the "underpowered for their own assertion" construction are
% retired by Amendment 2 and are quoted here as the superseded reading, never
% asserted. The current result is \AuditBelowThresholdAtMedian of
% \AuditAssessable at the registered uniform margin; see \S\ref{sec:audit:results}.
% Do NOT update these figures to rev-3 -- the point of the paragraph is that the
% rev-2 population change moved nothing, which is a statement about rev-2.
The headline verdicts did not move under the rev-2 population change. On the
own-margin reading then in force, since superseded by Amendment~2, the audit
stood at $K = 4$ of 12 determinate claims and $J = 5$ indeterminate, identical
to rev-1 under the rev-2 discordance imputation, with the uniform-2\,pp
secondary reading unchanged at 1 of 12. Those counts are quoted as the state of
the analysis at that date and are not current results: the eligible denominator
later changed with R10's exclusion, and the own-margin reading was withdrawn
outright. What survives is the finding that enlarging the atlas moved no
verdict. No new verdict computation was triggered. In the
certification tables no family entered or left, and \texttt{bbh}, \texttt{gpqa},
\texttt{ifeval}, \texttt{math}, \texttt{mmlu\_pro} and \texttt{musr} are
unchanged in both cell count and required $n$; the families that moved are
\texttt{mmlu} (798 $\to$ 1{,}311 cells, required $n$ 2{,}123 $\to$ 2{,}164),
\texttt{gsm8k} (11 $\to$ 24, 750 $\to$ 1{,}184), \texttt{winogrande}
(15 $\to$ 23, 1{,}879 $\to$ 1{,}416), \texttt{arc\_challenge} (8 $\to$ 17,
1{,}211 $\to$ 1{,}218), \texttt{hellaswag} (14 $\to$ 23, 688 $\to$ 695) and the
pooled row (1{,}155 $\to$ 1{,}707, 1{,}739 $\to$ 1{,}855).

That the audit's headline survived a 44\% enlargement of the atlas population is
worth stating plainly, and it is also worth stating that we did not know it
would when the repair was authorised.
